\VerdierPhysicalPageAnchor{0019}%
diagram commutes:

\VerdierSemanticAnchor{intro-1-resolution-comparison-diagram}
\[
\begin{tikzcd}[column sep=large,row sep=large]
X \arrow[r,"\varepsilon'"] \arrow[dr,"\varepsilon"']
  & \mathcal R(X) \arrow[d,"\rho"] \\
  & I^{\bullet}(X)
\end{tikzcd}
\]

\VerdierSemanticAnchor{intro-1-f-acyclic-morphism}
Moreover, given a morphism of $F$-acyclic resolutions:

\VerdierSemanticAnchor{intro-1-f-acyclic-morphism-diagram}
\[
\begin{tikzcd}[column sep=large,row sep=large]
X \arrow[r,"\varepsilon_1"] \arrow[dr,"\varepsilon_2"']
  & \mathcal R_1(X) \arrow[d,"\rho'"] \\
  & \mathcal R_2(X)
\end{tikzcd}
\qquad\text{,}
\]
the morphism of complexes of objects of $\mathcal B$:

\VerdierSemanticAnchor{intro-1-f-rho-prime}
\[
F(\rho') : F(\mathcal R_1(X)) \longrightarrow F(\mathcal R_2(X))
\]
is a quasi-isomorphism.

\VerdierSemanticAnchor{intro-1-inductive-system}
Observe that here one obtains in $\mathcal B$ a direct system of complexes and
quasi-isomorphisms having a unique maximal element up to homotopy: the image
under $F$ of an injective resolution of $X$. This phenomenon no longer occurs
in the second example, which we now discuss.

\VerdierSemanticAnchor{intro-1-left-derived-tensor}
The method of defining derived functors by resolutions sometimes generalizes
to the case where $\mathcal A$ does not necessarily have enough injective or
projective objects. Consider the case where $\mathcal A$ is the category of
sheaves of abelian groups on a topological space $E$. This category does not,
in general, have enough projective objects. We nevertheless propose to define
the \emph{left derived functors} of the right exact functor from $\mathcal A$
to $\mathcal A$ that associates to each sheaf $G$ the sheaf
$F\otimes_{\mathbb Z}G$ ($F$ a fixed sheaf). Naturally, we require this
definition to yield functors having the usual properties of left derived
functors (behavior with respect to exact sequences and a certain universal
property) and to recover, in the known cases (for example when $E$ is reduced
to a point), the usual torsion products.

\VerdierSemanticAnchor{intro-1-left-resolution}
One can proceed as follows. Every sheaf $G$ admits a left resolution
$L^{\bullet}\to G$ of finite length (in fact, of length 2) by
